\begin{vidu} %[2P4-6.1K][Loại 2]
Hai dây dẫn thẳng, dài song song và cách nhau 10 cm trong chân không, dòng điện trong hai dây cùng chiều có cường độ lần lượt là 2 A và 5 A. Lực từ tác dụng lên 20 cm chiều dài của mỗi dây có độ lớn
\choice
{lực đẩy có độ lớn $4.10^{-6}$ N}
{\True lực hút có độ lớn $4.10^{-6}$ N}
{lực đẩy có độ lớn $2.10^{-6}\ N$}
{lực hút có độ lớn $2.10^{-6}\ N$}
\loigiai{Hai dòng điện thẳng dài song song cùng chiều nên chúng hút nhau với lực có độ lớn $F = {2.10^{- 7}}\dfrac{{I_1}{I_2}\ell}{r} = {2.10^{- 7}}\dfrac{2.5.0,2}{0,1} = {4.10^{- 6}}$ N}
\haidong{2}
\end{vidu}
\begin{vidu} %[2P4-6.1B][Loại 3]
	Hai dây dẫn thẳng, dài song song đặt trong không khí. Dòng điện chạy trong hai dây có cùng cường độ 1 A. Lực từ tác dụng lên mỗi mét chiều dài có độ lớn là $10^{-7}$ N. Khoảng cách giữa hai dây là
	\choice
	{100 cm}
	{20 cm}
	{\True 200 cm}
	{10 cm}
	\loigiai{Độ lớn lực tác dụng lên mỗi mét chiều dài dây dẫn: $F = {2.10^{- 7}}\dfrac{{{I_1}{I_2}}}{r}$ $\Rightarrow$ khoảng cách giữa hai dây $r = {2.10^{- 7}}\dfrac{{{I_1}{I_2}}}{F} = {2.10^{- 7}}\dfrac{1.1}{{{{10}^{- 7}}}} = 2\ m = 200\ cm$}
\haidong{2}
\end{vidu}
\begin{vidu} %[2P4-6.1K][Loại 4]
	Hai dây dẫn thẳng dài vô hạn, song song đặt trong không khí cách nhau 1 cm, cường độ dòng điện chạy trong hai dây dẫn bằng nhau. Lực từ tác dụng lên mỗi mét chiều dài của mỗi dây là 2. $10^{-5}$ N. Cường độ dòng điện trong hai dây là
	\choice
	{0,1 A}
	{0,2 A}
	{\True 1 A}
	{2 A}
	\loigiai{Lực từ tác dụng lên mỗi mét chiều dài của mỗi dây là: $F={{2. 10}^{-7}}\dfrac{{I_1}{I_2}}{r}\Leftrightarrow {{2. 10}^{-5}}={{2. 10}^{-7}}\dfrac{{{I}^{2}}}{0,01}\Rightarrow I=1$ A}
\haidong{2}
\end{vidu}
\begin{vidu} %[2P4-6.1K][Loại 4]
	Có hai dây dẫn song song cách nhau 2 cm. Dòng điện chạy qua hai dây lần lượt là $I_1$ ,$I_2$ và $I_2 = 2,5I_1$. Khi lực từ tương tác lên mỗi mét chiều dài của mỗi dây là $10^{-4}$ N thì giá trị cường độ dòng điện qua mỗi dây là
	\choice
	{$I_1 = 5$ A; $I_2 = 2$ A}
	{$I_1 = 0,2$ A; $I_2 = 0,5$ A}
	{\True $I_1 = 2$ A; $I_2 = 5$ A}
	{$I_1 = 4$ A; $I_2 = 10$ A}
	\loigiai{$F={{2. 10}^{-7}}\dfrac{{I_1}{I_2}}{r} \Rightarrow I_1I_2 = 10$ mà $I_2 = 2,5I_1 \Rightarrow I_1 = 2$ A và $I_2 = 5$ A}
\haidong{4}
\end{vidu}
